\documentclass[final]{article}


\usepackage{aps360}
\usepackage[utf8]{inputenc} % allow utf-8 input
\usepackage[T1]{fontenc}    % use 8-bit T1 fonts
\usepackage{hyperref}       % hyperlinks
\usepackage{url}            % simple URL typesetting
\usepackage{booktabs}       % professional-quality tables
\usepackage{amsmath}        % math notation
\usepackage{amsfonts}       % blackboard math symbols
\usepackage{nicefrac}       % compact symbols for 1/2, etc.
\usepackage{microtype}      % microtypography
\usepackage{xcolor}         % colors
 
\title{Human Pose Estimation for Ergonomic Risk Scoring}

\author{%
  First \& Last Name \\
  Student Number, email\\
  Link to Google Colab or GitHub: TBD\\
}

\begin{document}

\maketitle

\vspace{-0.5in}

\section*{Introduction}
This project will build a real-time ergonomic risk classifier from webcam video. The motivating use case is human-robot collaboration: before a robot moves near a person, it should recognize risky postures such as trunk bending, neck strain, or overreaching. The final demo will overlay a skeleton, risk label, and confidence score on live video.

The task is four-class classification: \textit{neutral}, \textit{trunk\_flexion}, \textit{neck\_strain}, and \textit{overreach}. A pretrained pose estimator, such as MoveNet or MediaPipe Pose, will convert each frame into body keypoints. My trainable model will be a PyTorch multilayer perceptron (MLP) that maps normalized keypoints to four logits. Deep learning is useful because posture labels depend on several noisy joints, body sizes, and camera viewpoints. Success will be measured with accuracy, macro-F1, per-class precision/recall, a confusion matrix, and live inference speed.

\section*{Illustration}
\begin{figure}[!h]
  \centering
  \fbox{
  \begin{minipage}{0.92\linewidth}
  \centering
  \textbf{Webcam frame} $\rightarrow$
  \textbf{Pose estimator} $\rightarrow$
  \textbf{17 keypoints} $(x,y,c)$ $\rightarrow$
  \textbf{normalization} $\rightarrow$
  \textbf{MLP classifier} $\rightarrow$
  \textbf{risk label + overlay}
  \end{minipage}
  }
  \caption{The pretrained pose model extracts the skeleton; the trainable MLP predicts the ergonomic risk class.}
  \label{fig:pipeline}
\end{figure}

\section*{Background \& Related Work}

\section*{Data Processing}
The primary dataset will be MPOSE2021, a pose-based human action dataset with single-person keypoint sequences \citep{mazzia2021action,pic4ser2021mpose}. It is suitable because it provides skeleton data directly, letting the project focus on ergonomic classification rather than training a pose detector. I will map its keypoint order to the MoveNet/MediaPipe order used at webcam inference time.

MPOSE2021 has action labels, so I will repurpose it by generating ergonomic labels from joint geometry. For each pose, I will compute angles from three keypoints:
\[
\theta = \cos^{-1}\left(\frac{(A-B)\cdot(C-B)}{\|A-B\|\|C-B\|}\right),
\]
where $B$ is the measured joint. Simplified RULA/REBA thresholds \citep{mcatamney1993rula,hignett2000reba} will assign the four labels. If multiple risks occur, the largest threshold violation will decide the class.

I will remove samples with missing landmarks, filter low-confidence keypoints, center each skeleton at the mid-hip, and scale by torso length or shoulder width. Each frame will use $17 \times 3 = 51$ input values. Data will be split into train/validation/test sets by sequence or subject where possible. I will also record a small manually labelled webcam test set for the four classes.

\section*{Architecture}
The model is a small PyTorch MLP trained on 51-dimensional keypoint vectors:

\[
\text{Input}(51) \rightarrow \text{Linear}(128) \rightarrow \text{ReLU} \rightarrow
\text{Dropout} \rightarrow \text{Linear}(64) \rightarrow \text{ReLU} \rightarrow
\text{Linear}(4).
\]

The final layer outputs four class logits. I will train with cross-entropy loss and Adam, then tune learning rate, dropout, hidden size, and class weighting on the validation set. In the demo, OpenCV will draw the skeleton and predicted class on each webcam frame.

\section*{Baseline Model}
The baseline will be a rule-based classifier using the same normalized keypoints. It will compute neck, trunk, and upper-arm angles, then apply simplified RULA/REBA thresholds directly. If no threshold is violated, it predicts \textit{neutral}; otherwise it predicts the most severe violated class. This baseline is interpretable but sensitive to noisy keypoints. The MLP should improve robustness by learning joint patterns from data.

\section*{Ethical Considerations}

\newpage
\section*{References}
\begin{thebibliography}{9}

\bibitem[Bazarevsky et~al.(2020)]{bazarevsky2020blazepose}
Bazarevsky, V., Grishchenko, I., Raveendran, K., Zhu, T., Zhang, F., and Grundmann, M. (2020).
\newblock BlazePose: On-device real-time body pose tracking.
\newblock \textit{arXiv preprint arXiv:2006.10204}.

\bibitem[Hignett and McAtamney(2000)]{hignett2000reba}
Hignett, S. and McAtamney, L. (2000).
\newblock Rapid entire body assessment (REBA).
\newblock \textit{Applied Ergonomics}, 31(2), 201--205.

\bibitem[Mazzia et~al.(2021)]{mazzia2021action}
Mazzia, V., Angarano, S., Salvetti, F., Angelini, F., and Chiaberge, M. (2021).
\newblock Action Transformer: A self-attention model for short-time pose-based human action recognition.
\newblock \textit{Pattern Recognition}, 124, 108487.

\bibitem[McAtamney and Corlett(1993)]{mcatamney1993rula}
McAtamney, L. and Corlett, E. N. (1993).
\newblock RULA: A survey method for the investigation of work-related upper limb disorders.
\newblock \textit{Applied Ergonomics}, 24(2), 91--99.

\bibitem[PIC4SeR(2021)]{pic4ser2021mpose}
PIC4SeR. (2021).
\newblock MPOSE2021\_Dataset: A dataset for short-time human action recognition.
\newblock \url{https://github.com/PIC4SeR/MPOSE2021_Dataset}.

\bibitem[TensorFlow Hub(2024)]{tensorflowmovenet}
TensorFlow Hub. (2024).
\newblock MoveNet: Ultra fast and accurate pose detection model.
\newblock \url{https://www.tensorflow.org/hub/tutorials/movenet}.

\end{thebibliography}


\end{document}
